% !TEX root = ../tate.tex

\section{Preliminaries}

In this section we explain the properties of the $R$-tensor $\otimes_R$ and the ring dual $\Hom_R(-,R)$ of  $R$-modules. Both constructions can be generalised to chain complexes of $R$-modules by applying them degreewise.

\subsection{The $R$-tensor}
Let $R$ be a ring, $M$ a right $R$-module and $N$ a left $R$-module. The $R$-tensor $\otimes_R$ is defined by 
\[
M\otimes_R N=(M\otimes N)/(mr\otimes n)\sim (m\otimes rn).
\]
Note that as it is, $M\otimes_R N$ is just an abelian group, with no module structure.

We can think of  $\otimes_R$ as monoidal product in a category of bimodules, with objects triplets $(R, S, N)$ where $R$ and $S$ are rings and $N$ is an $(R,S)$-bimodule. Thus, for objects $(T, R, M)$ and $(R, S, N)$, their monoidal product is the object $(T,S,M\otimes_R N)$, so that $M\otimes_R N$ is a $(T, S)$-bimodule with left $T$-action given by $t(m\otimes n)=tm\otimes n$ and right $S$-action $(m\otimes n)s=m\otimes n s$.

If instead of a right $R$-module, $M$ is a left $R$-module, there are conditions under wich we can make sense of $\otimes_R$. We consider two situations for this.

\begin{example} %Note that we dont use this example at any point
If $R$ is a commutative ring, any left $R$-module can be turned into a bimodule by setting the right $R$-action to be the same as the left $R$-action, so that $m\cdot r:=rm$. Therefore
\[
M\otimes_R N=(M\otimes N)/(rm\otimes n)\sim (m\otimes rn).
\]
This is the monoidal product of the closed symmetric monoidal category $R$-Mod, which can be seen as the restriction of the category of bimodules to the full subcategory whose objects are triplets $(R,R, M)$ in which the left and right $R$-actions on $M$ coincide. Thus, $M\otimes_R N$ is again an $R$-module with $R$-action $r(m\otimes n)=rm\otimes n$.

The category $R$-Mod has an internal hom, $\Hom_R$, with $R$-module structure given by  $(r\cdot f)(m) = f(r m) = r f(m)$. Furthermore, for $R$-modules $M$, $N$, and $P$, we have the hom-tensor adjunction
\[
\Hom_R(M\otimes_R N, P)\cong \Hom_R(M,\Hom_R(N,P)).
\]
\end{example}
\begin{example} If $R=\k[G]$, for $\k$ a field and $G$ a group, we can turn a left $\k[G]$-module $M$ into a right $\k[G]$-module via the group anti-automorphism $g\mapsto g^{-1}$, by defining the right action to be (the $\k$-linear extension of) $m\cdot g:=g^{-1}\cdot m$. Note that this does not make $M$ into a bimodule unless $G$ is abelian because $(g\cdot m)\cdot h=g\cdot (m\cdot h)$ is only satisfied if $h^{-1}g m=g h^{-1} m$ for any $g,h\in G$.

We will denote the corresponding tensor by $\otimes_G$, to stress that we are using this specific right $\k[G]$-module structure on $M$, and is thus given by
\[
M\otimes_G N:=(M\otimes N)/(g^{-1}m\otimes n)\sim (m\otimes gn).
\]
Observe that replacing $m$ by $gm$ in the quotient relation gives $(m\otimes n) \sim (gm\otimes gn)$, which means that 
\[
M\otimes_G N = (M\otimes N)_G:= \k\otimes_{\k[G]}(M\otimes N),
\]
the coinvariants of $M\otimes N$ with respect to the diagonal left $G$-action $g\cdot (m\otimes n)= gm\otimes gn$. 

There is also a diagonal action on $\Hom(M,N)$ given by $(g\cdot f)(m)=g f(g^{-1}m)$, which gives the equivalence between $G$-equivariant maps and invariants
\[
\Hom_G(M,N)=(\Hom(M,N))^G:=\Hom_{\k[G]}(\k,\Hom(M,N)).
\]
Note that $\Hom_G(M,N)=\Hom_{\k[G]}(M,N)$ by definition.
%In particular, the linear dual $M^*:=\Hom_\k(M,\k)$ has a $\k[G]$-module structure given by $(g\cdot f)(m)=f(g^{-1}m)$.

Let $A$ be a $\k$-module (or rather, a $\k$-vector space), which we can regard as a $\k[G]$-module with the trivial action. We have the following identities regarding the functors of coinvariants and invariants:
\[
\Hom (M_G,A)=\Hom_G(M,A)\quad\text{and}\quad \Hom_G(A,M)=\Hom(A,M^G).
\]
\end{example}
\subsection{The ring dual}
The ring dual $M^\vee$ of a left $R$-module $M$ is the right $R$-module
\[
M^\vee=\Hom_{R}(M,R)
\]
 with action $(f\cdot r)(m)=(f(m))r$.
If $M$ is a finitely generated projective left $R$-module, then $M^\vee$ is a finitely generated projective right $R$-module, and for any left $R$-module $N$ we have 
\[
\varphi \colon M^\vee\otimes_{R}N\cong \Hom_{R}(M,N),
\]
given by $\varphi(f\otimes n)(m)=f(m)n$, since $f(m)\in R$.

\begin{example}
Let $R=\k[G]$ for $\k$ a field and $G$ a group, and consider the left $R$-module $W_0:=R\{e_0\}$ with the regular action $g\cdot e_0=ge_0$. We can also think of $W_0$ as $\k\{ge_0 \,|\, g\in G\}$, with action permuting the $\k$-basis. The dual $W_0^\vee$ is the right $R$-module
\[
W_0^\vee=\Hom_R(W_0,R)=\{e_0^\vee\}R=\k\{(ge_0)^\vee\, |\, g\in G\},
\]
where $(ge_0)^\vee\colon W_0\to R$ maps $R$-linearly $ge_0$ to $1_R$. The right action is
\[
e_0^\vee\cdot g=(g^{-1}e_0)^\vee,
\]
since $(e_0^\vee\cdot g)(he_0)=(e_0^\vee(he_0))g=hg=(g^{-1}e_0)^\vee(hgg^{-1}e_0)=(g^{-1}e_0)(he_0)$ for any $h\in G$.

Recall that we can define a right $R$-module structure on $W_0$ by $e_0\cdot g=g^{-1}e_0$. If $G$ is abelian, then $W_0$ is a bimodule and we can also define the left $R$-action on $W_0^\vee$ 
\[
(g\cdot e_0^\vee) (he_0) := e_0^\vee(he_0\cdot g)= (ge_0)^\vee(he_0),
\]
where the second equality holds by abelianity of $G$ (or the bimodule structure of $W_0$), for 
\[
e_0^\vee(he_0\cdot g)=g^{-1}h=hg^{-1}=(ge_0)^\vee(hg^{-1}ge_0).
\]
Note that the bimodule structure is also needed to ensure that $g\cdot e_0^\vee\in \Hom_R(M,R)$,
\[
(g\cdot e_0^\vee)(he_0)=e_0^\vee(he_0\cdot g)=e_0^\vee(h(e_0\cdot g))=h(g\cdot e_0^\vee).
\]
Observe that we could have defined the right action on $W_0^\vee$  as the action given by acting on the left with the group inverse
\[
e_0^\vee\cdot g:=g^{-1}\cdot e_0^\vee=(g^{-1}e_0)^\vee,
\]
which means that it corresponds to the right action for the $G$-tensor $W_0^\vee\otimes_G -$.
\end{example}

\section{Main}

\subsection{The Tate construction}

Let \(\k\) be a field, let \(G\) be a finite group, and set \(R = \k[G]\).
We denote by \(\Ch(R)\) the category of unbounded chain complexes of left \(R\)-modules equipped with the projective model structure, weak equivalences are quasi-isomorphisms and fibrations are degreewise surjections.
This category is pointed, with zero object the complex~\(0\).
For a map \(f \colon X \to Y\) in \(\Ch(R)\), the \defn{cofibre} \(\cofib(f)\) and \defn{fibre} \(\fib(f)\) are defined by the pushout and pullback squares
\[
\begin{tikzcd}
	X \ar[r, "f"] \ar[d] & Y \ar[d] \\
	0 \ar[r] & \cofib(f)
\end{tikzcd}
\qquad \text{and} \qquad
\begin{tikzcd}
	\fib(f) \ar[r] \ar[d] & 0 \ar[d] \\
	X \ar[r, "f"] & Y
\end{tikzcd}
\]
respectively.

We denote the homotopy category of \(\Ch(R)\) by \(\sD(R)\), its localization at quasi-isomorphisms.
The suspension functor \([1]\) is induced by the shift of complexes and is an autoequivalence of \(\sD(R)\).
Moreover, homotopy pushouts and homotopy pullbacks coincide in \(\sD(R)\), and in particular \(\hofib(f) \simeq \hocofib(f)[-1]\) for any map \(f\) in \(\sD(R)\).
We refer to sequences \(X \to Y \to \hocofib(f)\) and \(\hofib(f) \to X \to Y\) as \defn{cofibre sequences} and \defn{fibre sequences}, respectively; by the above, they carry equivalent information.

The functors
\[
(-)_G = \k \ \otimes_G \, -
\qquad \text{and} \qquad
(-)^G = \Hom_G(\k, -)
\]
from \(\Ch(R)\) to \(\Ch(\k)\) induce the \defn{derived coinvariants} and \defn{derived invariants} functors from \(\sD(R)\) to \(\sD(\k)\).
Explicitly, for \(M \in \sD(R)\) these are given by
\[
M_{hG} \coloneqq W \otimes_G M
\qquad \text{and} \qquad
M^{hG} \coloneqq \Hom_G(W, M),
\]
where \(W \xrightarrow{\sim} \k\) is a cofibrant resolution.
The homology groups of \(M_{hG}\) and \(M^{hG}\) recover group homology \(H_*(G; M)\) and group cohomology \(H^*(G; M)\), respectively.

There is a natural transformation
\[
\Nm \colon (-)_{hG} \to (-)^{hG}
\]
called the \defn{norm map}, which derives the classical map \(M_G \to M^G\) given by the action of the norm element \(N = \sum_{g \in G} g \in R\).
Concretely, since \(\Hom_R(\k, R) \cong \k\), the \(R\)-dual \(W^\vee = \Hom_R(W, R)\) is itself a cofibrant resolution of \(\k\); any chain homotopy equivalence \(W \xrightarrow{\sim} W^\vee\) then induces the norm map via
\[
W \otimes_R M \longrightarrow W^\vee \otimes_R M \cong \Hom_R(W, M),
\]
where the final isomorphism holds because \(W\) consists of finitely generated projectives.

The \defn{Tate construction} is the functor \((-)^{tG} \defeq \hocofib(\Nm) \colon \sD(R) \to \sD(\k)\).
The homology of \(M^{tG}\) agrees with the \defn{Tate cohomology groups} \(\widehat{H}^*(G; M)\) when \(M\) is concentrated in degree 0.
The resulting cofibre sequence
\[
M_{hG} \xrightarrow{\Nm} M^{hG} \longrightarrow M^{tG}
\]
in \(\sD(\k)\) relates group homology, group cohomology, and Tate cohomology. %It is cohomology indeed.

\subsection{Explicit models for \(\cyc_p\) over \(\Fp\)}

Let \(p\) be a prime and \(\cyc_p = \angles{\rho \mid \rho^p}\) be the cyclic group of order \(p\).
We work over the field \(\Fp\) with \(p\) elements and denote \(R = \Fp[\cyc_p]\).
Set \(T = \rho - 1\) and \(N = 1 + \rho + \cdots + \rho^{p-1}\) in \(R\), as introduceed before.
A cofibrant resolution \(W(p) \xrightarrow{\sim} \Fp\) is given by the complex
\[
W(p) = R\set{e_0} \xla{T} R\set{e_1} \xla{N} R\set{e_2} \xla{T} R\set{e_3} \xla{N} \cdots,
\]
where each \(R\set{e_i}\) is a free \(R\)-module of rank one concentrated in homological degree \(i \geq 0\).
Since each \(R\set{e_i} \cong R\) is finitely generated and free, the \(R\)-dual \(W^\vee(p) = \Hom_R(W, R)\) is the complex
\[
W^\vee(p) = \cdots \xla{T} \set{e_2^\vee}R \xla{N} \set{e_1^\vee}R \xla{T} \set{e_0^\vee}R.%The R's could go in the left if we are considering the bimodule structure of Example 3, although the natural action is the right R action.
\]

Let \(N_{W(p)} \colon W(p) \xrightarrow{\sim} W^\vee(p)\) be the map that is \(0\) in positive and negative degrees and sends \(e_0\) to \(N e_0^\vee\).
Its mapping cone is isomorphic to
\[
\cone(N_{W(p)}) \, \cong \, \widehat{W}(p) = \cdots \xla{N} R\set{e_{-1}} \xla{T} R\set{e_0} \xla{N} R\set{e_1} \xla{T} R\set{e_2} \xla{N} \cdots
\]
with \(e_i\) in degree \(i\).

Consider the cofibre sequence
\[
W(p) \to W^\vee(p) \to \widehat W(p).
\]
For any chain complex \(M\), applying the functor \(- \ot_R M\) to it gives a cofibre sequence representing
\begin{equation}
M_{h\cyc_p} \xra{\Nm} M^{h\cyc_p} \to M^{t\cyc_p}.
\label{CofibreM}
\end{equation}

%\subsection{Diagonals}
%
%Let $C$ be a chain complex. An \defn{unstable $p$-cyclic diagonal} on  $C$, or \defn{unstable $p$-diagonal} for short, is a $\cyc_p$-equivariant chain map
%	\[
%	\mu \colon W(p) \otimes C \to C^{\otimes p},
%	\]
%where the $\cyc_p$-action on the domain is given by multiplication on the first factor and on the codomain consists of cyclically permuting the factors.
%
%The fibre of $N_{W(p)}$ is isomorphic to the desuspension of $\widehat W(p)$, which we denote by
%\[
%W^{\mathrm st}(p) \, := \, \widehat{W}(p)[-1] = \cdots \xla{T} R\set{e_{-1}} \xla{N} R\set{e_0} \xla{T} R\set{e_1} \xla{N} R\set{e_2} \xla{T} \cdots
%\]
%with $e_i$ in degree $i$. An \defn{stable $p$-cyclic diagonal} on  $C$, or simply a\defn{stable $p$-diagonal}, is a $\cyc_p$-equivariant chain map
%	\[
%	\mu \colon W^{\mathrm st}(p) \otimes C \to C^{\otimes p}.
%	\]
%	
%We can recover both notions of diagonals from the cofibre sequence \eqref{CofibreM} by setting $M=\Hom(C,C^{\otimes p})$, which gives
%\[
%W(p)\otimes_R \Hom(C,C^{\otimes p}) \xra{\Nm} W^\vee(p)\otimes_R \Hom(C,C^{\otimes p})\to \widehat W(p)\otimes_R \Hom(C,C^{\otimes p}).
%\]
%Since $W$ consists of finitely generated projectives, we can express it in terms of homs
%\[
%\Hom_R(W^\vee(p), \Hom(C,C^{\otimes p})) \xra{\Nm} \Hom_R(W(p), \Hom(C,C^{\otimes p}))\to \Hom_R(W^{\mathrm st}(p), \Hom(C,C^{\otimes p})),
%\]
%and by hom-tensor adjunction %%% This is not clear yet
%\[
%\Hom(W^\vee(p)\otimes_R C,C^{\otimes p}) \xra{\Nm} \Hom(W(p)\otimes_R C,C^{\otimes p}))\to \Hom(W^{\mathrm st}(p)\otimes_RC,C^{\otimes p})).
%\]
%Thus, unstable and stable $p$-cyclic diagonals  are, respectively,  $0$-cycles in the second and third chain complexes of the cofibre sequence.
\subsection{The algebraic Tate diagonal} Let $M= C^{\otimes p}$. Then, a model for its tate construction is
\[
(C^{\otimes p})^{t C_p}\cong \widehat W(p)\otimes_R C^{\otimes p}.
\]
The no-go theorem from Nikoulaus and Scholze \cite{??} states that there is no natural chain map 
\[
C\to (C^\otimes p)^{tC_p}
\]
that induces a non trivial map in homology. However, we can construct the map 
\[
\Delta^t\colon H_*(C)\to H_*((C^\otimes p)^{tC_p})=H_*(\widehat W\otimes_R C^{\otimes p})
\]
that sends a homology class $[x]$ in degree $n$ to $[e_{n-pn}\otimes_R (x\otimes \overset{_{(p)}}{\cdots}\otimes x)]$. We call $\Delta^t$ the \defn{algebraic Tate diagonal}. \TBW
...